
\sgaStableTarget{sga6.e10.prop.6-2-2}\textbf{Proposition 6.2.2.} \uline{For every positive section} @@M0@@ \uline{of} @@M1@@\uline{, the étale sheaf of groups} @@M2@@ \uline{of commutators of} @@M3@@ \uline{is represented by}
@@M4@@.

\sgaStableTarget{sga6.e10.cor.6-2-3}\textbf{Corollary 6.2.3.} \uline{Every representation}
\sgaStableTarget{sga6.e10.eq.6-2-4}@@M5@@
\uline{is of the form}
\sgaStableTarget{sga6.e10.eq.6-2-5}@@M6@@
\uline{where} @@M7@@, @@M8@@ \uline{and} @@M9@@ \uline{are sections of} @@M10@@ \uline{uniquely determined by} @@M11@@.

\sgaStableTarget{sga6.e10.cor.6-2-6}\textbf{Corollary 6.2.6.} \uline{With the notation of} 6.1. \uline{one has}
\sgaStableTarget{sga6.e10.eq.6-2-6}@@M12@@
\uline{where} @@M13@@ \uline{are sections of} @@M14@@ \uline{uniquely determined by} @@M15@@.

\sgaStableTarget{sga6.e10.rem.6-2-7}\textbf{Remark 6.2.7.} One can determine the exponents @@M16@@ in (6.2.5) and (6.2.6) by successively substituting in (6.2.5)
@@M17@@, @@M18@@, for @@M19@@ @@M20@@, whence
@@M21@@
For example, one thus deduces from (6.2.6)
\sgaStableTarget{sga6.e10.eq.6-2-8}\sgaStableTarget{sga6.e10.eq.6-2-9}@@M22@@
and
\sgaStableTarget{sga6.e10.eq.6-2-10}@@M23@@
etc.
